۴۰۸  
بحث خود را کوتاه کردیم زمانی که میستری با گام‌هایی محکم وارد آشپزخانه شد.  
«تمام شد،» گفت. «با نامزد آنیا در ملز صحبت طولانی داشتم. به او گفتم که برای جبران کارهایش با او دیر شده است. حالا آنیا دوست‌دختر من است و ما عاشق هم هستیم. این بهترین موفقیتی است که در تاریخ روش میستری اتفاق افتاده است.»  
مارینا به من نگاهی آشنا انداخت. مادر میستری دست‌هایش را روی سینه‌اش جمع کرد و به خود لبخندی زد.  

من لیسا را ساعت ۸ شب برداشتم و او را به رستورانی ژاپنی به نام "کتانا" بردم. او ضبط صوتی را با ضربه روی اپن آشپزخانه گذاشت. «کل مکالمه را ضبط کردم،» گفت. «یکی از سخت‌ترین شام‌هایی بود که داشتم. ما بقدری وقت با هم گذرانده بودیم که حتی دیگر موزه‌ای نداشتم. مجبور بودم خودم باشم.»  
«نه،» به او گفتم. از این همه درام خسته شده بودم.  
علاوه بر این، من قراری با لیسا داشتم که باید به آن عمل می‌کردم. «چیزی هست که مدتی است می‌خواهم بپرسم،» گفتم در حالی که لامپ‌های حرارتی کنار ایوان رستوران فرق سرمان را داغ کرده بود و ساکه معده‌هایمان را گرم می‌کرد. این سؤال هفته‌ها بود که باعث بی‌خوابی من شده بود. «بعد از آتلانتا چه بر سرت آمد؟ ما برنامه داشتیم و تو آن‌ها را نقض کردی.»  
«تو در تلفن بی‌ادب بودی،» او گفت. «و من فکر نمی‌کردم ما برنامه‌های قطعی هم داشتیم.» بنابراین این نسخه او از نظریه طناب گربه بود، مجازات کردن برای رفتار بد.  

«من داشتم شوخ طبعی متکبرانه می‌کردم. می‌خواستم تو را ببینم.»  
«هر چه. تو بی‌ادب بودی. تو زیادی خودت را خونسرد و بی‌خیال و جدا از دیگران نشان دادی که خسته‌کننده بود. به خودم گفتم، ‘می‌توانم هر شخصی را به دست بیاورم و ناگهان این آدم در نقش آقای کول است؟’»  
هنگام صحبت کردن، سعی می‌کردم بفهمم چرا این دختر را این‌قدر دوست دارم، چرا بعد از آشنایی با این همه آدم، او برای من تبدیل به مشتاقیت شده است. بخش بدبینانه‌ای از من می‌گفت که من به سادگی در حال عاشق شدن به معادل مؤنث تاکتیک‌هایی هستم که استفاده می‌کنیم. راز این‌که چگونه کسی فکر کند که عاشق تو است، این است که افکارت را به خود مشغول کند، و این همان کاری است که لیسا با من کرده است. او مرا نادیده گرفته و از لحاظ جسمی مرا رد کرده بود در حالی که با تشویق کافی مرا به دنبال خود نگه داشته بود.  

از سوی دیگر، من کسی نبودم که به هر قیمتی به دنبال خواسته بروم. اگر زنی که برایم مهم نبود این‌قدر خود را دور از دسترس نشان می‌داد، مدت‌ها پیش منصرف می‌شدم. البته، ممکن است مشتاقیتم از یک رگه مردانه و زورآورانه ناشی شود که به طور تصادفی به عنوان یکی از عوارض جانبی تلاش‌های اجتماعی به دست آورده‌ام. لیسا به شدتی مستقل بود، کسی که من به او به جای نگاه به پایین، نگاه بالا می‌کردم. بنابراین شاید مرد غارنشین درونی‌ام فقط می‌خواست با او بخوابد و در نتیجه او را تسخیر کند.  

و همیشه این احتمال دور بود که او بتواند به قسمتی از من دست یابد که از همه، حتی خودم، مخفی نگه داشته‌ام. این قسمتی از من بود که می‌خواست دست از فکر کردن بردارد، دست از جست‌وجو بردارد، دست از نگرانی بردارد.